\section{Running jobs side by side}
\label{sec:interference}

\begin{question}
  \textbf{Can jobs run side by side without slowing each other, and does the background load move our
  numbers?} If four 13-node jobs run as fast together as alone, the later experiments can be packed in
  parallel. The fabric and filesystem are shared with other partitions, so the same run may also change
  speed during the window.
\end{question}

\paragraph{Interference (E1).} The Qwen best guess (HSDP within a node, 24k tokens per GPU) runs on 4
disjoint random sets of 13 nodes: each set alone, then the 4 together. This is repeated twice with new
random sets. If the slowdown together is below 2\,\%, runs of up to 26 nodes are packed in parallel for
the layout, scaling, precision and operations experiments. The median change together vs.\ alone is
\TriInterferencePct\,\%.

\begin{figure}[h]
  \centering\fig[0.72\linewidth]{interference}
  \caption{Tokens/s per GPU of each 13-node set (2 draws of 4 sets), alone (circles) and with the 3
    others running (squares). \syntheticnote{interference}}
  \label{fig:interference}
\end{figure}

\placeholder{whether any set slows down beyond the run-to-run spread; decision on packing.}

\paragraph{Reference drift (R).} The reference run (\cref{sec:method}) is repeated about every 3 hours.
Its spread bounds the differences between configurations that can be attributed to the configurations
themselves: the fastest reference run is \TriDriftPct\,\% faster than the slowest.

\begin{figure}[h]
  \centering\fig[0.72\linewidth]{reference_drift}
  \caption{Throughput of the reference run (Qwen, 4 fixed nodes, FSDP, 12{,}288 tokens per GPU) vs.\ time
    since the first reference run. Band: median $\pm1\,\%$; the linear trend over the window is written
    on top. \syntheticnote{reference_drift}}
  \label{fig:reference_drift}
\end{figure}

\placeholder{correlation with time of day or with other users' load.}
